% GENERATED FILE. Edit canonical lessons, then run npm run generate:books.
% canonical-source-hash: fnv1a64:4a16bd23fe0daa72
% canonical-lessons: KA-C238-taip-madu, KA-C238-kaleduhogu, KA-C238-oppu, KA-C238-nirakarisu, KA-C238-sampadisu

\chapter{To type, To get lost, To agree, To refuse, To earn}
\label{ch:ka-to-type-to-get-lost-to-agree-to-refuse-to-earn}
\hlchaptermodality{238}

\begin{chapteropening}
\textbf{By the end of this chapter:} \emph{I can say to type, to get lost, to agree, to refuse and to earn.}

Complete the last lesson of chapter 238: I can say to type, to get lost, to agree, to refuse and to earn.
\end{chapteropening}

\section[\texorpdfstring{\kn{ಟೈಪ್} \kn{ಮಾಡು} (ṭaip māḍu)}{ಟೈಪ್ ಮಾಡು (ṭaip māḍu)}]{\kn{ಟೈಪ್} \kn{ಮಾಡು} (ṭaip māḍu) — to type}

\label{lesson:KA-C238-taip-madu}



\begin{quote}
Before the new one: say the Kannada for to hit, then the Kannada for to try.
\end{quote}

\subsection*{You'll want to know: \kn{ಟೈಪ್} \kn{ಮಾಡು}}
\textbf{\kn{ಟೈಪ್} \kn{ಮಾಡು}} — \emph{ṭaip māḍu} — \textquotedblleft{}to type\textquotedblright{}.

\begin{grammarlens}[title={What you've built}]
One more word for this chapter.
\end{grammarlens}

\subsection*{Your turn}
\begin{itemize}
\raggedright
  \item \emph{Say it:} \emph{ṭaip māḍu}
  \item \emph{Say it:} \emph{ṭaip māḍu}, once more
  \item \emph{Say it:} say \emph{ṭaip māḍu}
  \item \emph{Recall:} say the Kannada for to hit, then the Kannada for to try, then say \emph{ṭaip māḍu} again
\end{itemize}

\begin{tcolorbox}[breakable,colback=teal!4,colframe=teal!35!black,title={Before you move on}]
What does \kn{ಟೈಪ್} \kn{ಮಾಡು} mean? (\textquotedblleft{}To type\textquotedblright{}.) Say it once more.
\end{tcolorbox}
\section[\texorpdfstring{\kn{ಕಳೆದುಹೋಗು} (kaḷeduhōgu)}{ಕಳೆದುಹೋಗು (kaḷeduhōgu)}]{\kn{ಕಳೆದುಹೋಗು} (kaḷeduhōgu) — to get lost}

\label{lesson:KA-C238-kaleduhogu}



\begin{quote}
Before the new one: say the Kannada for to try, then the Kannada for to type.
\end{quote}

\subsection*{You'll want to know: \kn{ಕಳೆದುಹೋಗು}}
\textbf{\kn{ಕಳೆದುಹೋಗು}} — \emph{kaḷeduhōgu} — \textquotedblleft{}to get lost\textquotedblright{}.

\begin{grammarlens}[title={What you've built}]
One more word for this chapter.
\end{grammarlens}

\subsection*{Your turn}
\begin{itemize}
\raggedright
  \item \emph{Say it:} \emph{kaḷeduhōgu}
  \item \emph{Say it:} \emph{kaḷeduhōgu}, once more
  \item \emph{Say it:} say \emph{kaḷeduhōgu}
  \item \emph{Recall:} say the Kannada for to try, then the Kannada for to type, then say \emph{kaḷeduhōgu} again
\end{itemize}

\begin{tcolorbox}[breakable,colback=teal!4,colframe=teal!35!black,title={Before you move on}]
What does \kn{ಕಳೆದುಹೋಗು} mean? (\textquotedblleft{}To get lost\textquotedblright{}.) Say it once more.
\end{tcolorbox}
\section[\texorpdfstring{\kn{ಒಪ್ಪು} (oppu)}{ಒಪ್ಪು (oppu)}]{\kn{ಒಪ್ಪು} (oppu) — to agree}

\label{lesson:KA-C238-oppu}



\begin{quote}
Before the new one: say the Kannada for to type, then the Kannada for to get lost.
\end{quote}

\subsection*{You'll want to know: \kn{ಒಪ್ಪು}}
\textbf{\kn{ಒಪ್ಪು}} — \emph{oppu} — \textquotedblleft{}to agree\textquotedblright{}.

\begin{grammarlens}[title={What you've built}]
One more word for this chapter.
\end{grammarlens}

\subsection*{Your turn}
\begin{itemize}
\raggedright
  \item \emph{Say it:} \emph{oppu}
  \item \emph{Say it:} \emph{oppu}, once more
  \item \emph{Say it:} say \emph{oppu}
  \item \emph{Recall:} say the Kannada for to type, then the Kannada for to get lost, then say \emph{oppu} again
\end{itemize}

\begin{tcolorbox}[breakable,colback=teal!4,colframe=teal!35!black,title={Before you move on}]
What does \kn{ಒಪ್ಪು} mean? (\textquotedblleft{}To agree\textquotedblright{}.) Say it once more.
\end{tcolorbox}
\section[\texorpdfstring{\kn{ನಿರಾಕರಿಸು} (nirākarisu)}{ನಿರಾಕರಿಸು (nirākarisu)}]{\kn{ನಿರಾಕರಿಸು} (nirākarisu) — to refuse}

\label{lesson:KA-C238-nirakarisu}



\begin{quote}
Before the new one: say the Kannada for to get lost, then the Kannada for to agree.
\end{quote}

\subsection*{You'll want to know: \kn{ನಿರಾಕರಿಸು}}
\textbf{\kn{ನಿರಾಕರಿಸು}} — \emph{nirākarisu} — \textquotedblleft{}to refuse\textquotedblright{}.

\begin{grammarlens}[title={What you've built}]
One more word for this chapter.
\end{grammarlens}

\subsection*{Your turn}
\begin{itemize}
\raggedright
  \item \emph{Say it:} \emph{nirākarisu}
  \item \emph{Say it:} \emph{nirākarisu}, once more
  \item \emph{Say it:} say \emph{nirākarisu}
  \item \emph{Recall:} say the Kannada for to get lost, then the Kannada for to agree, then say \emph{nirākarisu} again
\end{itemize}

\begin{tcolorbox}[breakable,colback=teal!4,colframe=teal!35!black,title={Before you move on}]
What does \kn{ನಿರಾಕರಿಸು} mean? (\textquotedblleft{}To refuse\textquotedblright{}.) Say it once more.
\end{tcolorbox}
\section[\texorpdfstring{\kn{ಸಂಪಾದಿಸು} (sampādisu)}{ಸಂಪಾದಿಸು (sampādisu)}]{\kn{ಸಂಪಾದಿಸು} (sampādisu) — to earn}

\label{lesson:KA-C238-sampadisu}



\begin{quote}
Before the new one: say the Kannada for to agree, then the Kannada for to refuse.
\end{quote}

\subsection*{You'll want to know: \kn{ಸಂಪಾದಿಸು}}
\textbf{\kn{ಸಂಪಾದಿಸು}} — \emph{sampādisu} — \textquotedblleft{}to earn\textquotedblright{}.

\begin{grammarlens}[title={What you've built}]
One more word for this chapter.
\end{grammarlens}

\subsection*{Your turn}
\begin{itemize}
\raggedright
  \item \emph{Say it:} \emph{sampādisu}
  \item \emph{Say it:} \emph{sampādisu}, once more
  \item \emph{Say it:} say \emph{sampādisu}
  \item \emph{Recall:} say the Kannada for to agree, then the Kannada for to refuse, then say \emph{sampādisu} again
\end{itemize}

\begin{tcolorbox}[breakable,colback=teal!4,colframe=teal!35!black,title={Before you move on}]
What does \kn{ಸಂಪಾದಿಸು} mean? (\textquotedblleft{}To earn\textquotedblright{}.) Say it once more.
\end{tcolorbox}
